% Section 4.2 — Mitigation Methods
% Paste under \subsection{Mitigation Methods} inside Section 4.

\subsection{Mitigation Methods}
\label{sec:mitigations}

We evaluate four published mitigation strategies alongside our own method.
All baselines operate \emph{downstream} of the preference data — either at the
reward-model training stage or the policy-optimization stage — and are therefore
blind to the missingness structure $\psi(a)$ that generates the bias.
Our method intervenes \emph{upstream}, at the preference-data stage, before
the reward model ever sees the data.
We deliberately tune each baseline to minimize round-one disparity (see
\S\ref{sec:tuning}), giving each the strongest possible chance of breaking the
amplification loop.

%% ─────────────────────────────────────────────────────────────────────────── %%
\paragraph{Baseline 1: Length normalization.}
Long responses tend to receive higher reward-model scores independent of
quality~\citep{singhal2023longisbet, park2024disentangling}, and groups whose
training pairs are disproportionately short may therefore receive systematically
lower scores.
Length normalization subtracts a length penalty from the Bradley--Terry
objective at training time.
For a preference pair $(x, y_w, y_l)$ the corrected reward is

\begin{equation}
    \tilde{r}(x, y) \;=\; r_\theta(x, y) \;-\; \beta \log |y|,
\label{eq:length_norm}
\end{equation}

where $|y|$ is the token count of response $y$ and $\beta > 0$ is a
hyperparameter.
We sweep $\beta \in \{0.01, 0.1, 1.0\}$ and select the value that minimizes
round-one disparity on a held-out validation set.
Length normalization is also applied inside the DPO objective by normalizing
per-token log-probabilities by sequence length before summing.

\emph{Why it is insufficient.}
The bias in our setting arises from the \emph{coverage} of demographic groups
in the training set, not from a systematic response-length difference between
groups.
Because Group~A and Group~B resumes are sampled from the same quality and
length distributions, length normalization has no mechanism to address the
underlying missingness gap and is expected to leave the amplification ratio
$\mathrm{AR}(r)$ unchanged after round one.

%% ─────────────────────────────────────────────────────────────────────────── %%
\paragraph{Baseline 2: Reward-model regularization.}
We regularize the reward head via two complementary penalties added to the
Bradley--Terry loss:

\begin{equation}
    \mathcal{L}_{\mathrm{reg}} \;=\;
    \mathcal{L}_{\mathrm{BT}}
    \;+\; \lambda_{\ell_2} \lVert W_r \rVert_F^2
    \;+\; \lambda_{\mathrm{cons}} \cdot
    \frac{1}{|G|}\sum_{g \in G} \operatorname{Var}_{(x,y)\sim\mathcal{D}_g}
    \bigl[r_\theta(x,y)\bigr],
\label{eq:rm_reg}
\end{equation}

where $W_r$ is the reward head weight matrix, $G = \{A, B\}$ denotes the
demographic groups, $\mathcal{D}_g$ is the subset of training pairs whose
chosen response belongs to group $g$, and $\operatorname{Var}[\cdot]$ is
estimated within each mini-batch.
The first term is a standard Frobenius-norm penalty (separate from the AdamW
weight decay, which applies to all parameters).
The second is a \emph{logit-consistency} term that discourages the reward head
from assigning high-variance scores within demographic strata, reducing the
signal-to-noise ratio of any spurious demographic shortcut.
We sweep $(\lambda_{\ell_2}, \lambda_{\mathrm{cons}}) \in \{10^{-4}, 10^{-3}\}^2$
and select the best round-one configuration.

\emph{Why it is insufficient.}
Regularization constrains the RM's capacity but does not change the
\emph{composition} of the training set.
After the RM has been trained on predominantly Group-A pairs, even a heavily
regularized head will have seen fewer Group-B examples per cell and will
therefore produce lower-confidence, less-calibrated scores for Group-B.
The loop re-uses these scores to generate round-$(r{+}1)$ preferences, so the
coverage gap compounds across rounds regardless of the regularization
applied in a single round.

%% ─────────────────────────────────────────────────────────────────────────── %%
\paragraph{Baseline 3: KL-constrained policy optimization.}
A standard remedy for reward overoptimization is to constrain how far the
policy $\pi_\theta$ may drift from the reference policy $\pi_{\mathrm{ref}}$:

\begin{equation}
    \max_{\pi_\theta}\;
    \mathbb{E}_{x \sim \mathcal{D},\, y \sim \pi_\theta(\cdot|x)}
    \bigl[r_\theta(x, y)\bigr]
    \;-\;
    \gamma \, \mathrm{KL}\!\bigl(\pi_\theta(\cdot|x) \,\|\, \pi_{\mathrm{ref}}(\cdot|x)\bigr).
\label{eq:kl_ppo}
\end{equation}

We implement this via RLOO~\citep{kool2019buy}, TRL's KL-constrained
online RL trainer, with $\gamma \in \{10^{-4}, 10^{-3}, 10^{-2}\}$
following~\citet{wolf2024tradeoffs}.
A larger $\gamma$ prevents the policy from exploiting reward artifacts but
also slows learning of genuine quality improvements.

\emph{Why it is insufficient.}
The KL constraint operates on the policy, not on the reward model or the
preference data.
The RM that scores policy outputs in each round is itself trained on a biased
sample; constraining how much the policy moves toward that biased RM's signal
reduces the \emph{rate} of bias amplification but does not reverse it.
As shown in \S\ref{sec:results}, $\mathrm{AR}(r)$ remains above 1.0 for all
values of $\gamma$ by round 5.

%% ─────────────────────────────────────────────────────────────────────────── %%
\paragraph{Baseline 4: Safe-RLHF fairness penalty.}
Safe-RLHF~\citep{dai2023safe} augments the reward with a safety signal that
penalizes high-cost behavior.
We adapt this framework to demographic fairness by replacing the safety cost
with the within-batch demographic parity gap:

\begin{equation}
    \tilde{r}_{\mathrm{safe}}(x, y) \;=\;
    r_\theta(x, y)
    \;-\;
    \lambda_{\mathrm{safe}} \cdot
    \max\!\Bigl(0,\;
    \overline{r}_A^{\,\mathrm{batch}} - \overline{r}_B^{\,\mathrm{batch}}
    \Bigr),
\label{eq:safe_rlhf}
\end{equation}

where $\overline{r}_g^{\,\mathrm{batch}}$ is the mean reward for group $g$ in
the current mini-batch, estimated from the demographic metadata attached to each
preference pair.
We sweep $\lambda_{\mathrm{safe}} \in \{0.1, 0.5, 1.0\}$.

\emph{Why it is insufficient.}
Equation~\eqref{eq:safe_rlhf} penalizes the policy for outputs that the
\emph{current-round RM} scores unequally across groups.
However, if the RM is already biased — because it was trained on a coverage-
unbalanced dataset — its scores are themselves a biased proxy for true quality.
Penalizing the policy based on a biased scorer is analogous to correcting a
thermometer reading with the same thermometer: the disparity signal the penalty
acts on is corrupted by the very bias we wish to remove.

%% ─────────────────────────────────────────────────────────────────────────── %%
\paragraph{Our method: Inverse propensity weighting (DPO$^+$).}
All four baselines share a common failure mode: they intervene \emph{after} the
biased preference data has been used to train the RM.
Our method instead reweights the preference pairs \emph{before} RM training,
blocking the backdoor path from demographic group $a$ through
missingness $m$ to RM score.

Recall from \S\ref{sec:setup} that pair $(x, y_w, y_l)$ is observed in the
training set with probability $\psi(a) = P(m{=}1 \mid a)$, where $a$ is the
demographic group of the candidate.
We estimate $\hat\psi(a)$ via a propensity model (oracle in the synthetic
setting; a logistic-regression classifier on prompt features in the real-world
setting) and reweight the Bradley--Terry loss:

\begin{equation}
    \mathcal{L}_{\mathrm{IPW}}(\theta) \;=\;
    -\sum_{i=1}^{N}
    \frac{1}{\hat\psi(a_i)}
    \log \sigma\!\bigl(r_\theta(x_i, y_i^w) - r_\theta(x_i, y_i^l)\bigr),
\label{eq:ipw}
\end{equation}

with propensities clipped to $[\varepsilon, 1{-}\varepsilon]$,
$\varepsilon = 0.05$, to prevent numerical blow-up.
Pairs from under-observed groups (low $\hat\psi$) receive higher weight,
compensating for their lower representation in the training set.
The same weights are passed to the DPO trainer when using the DPO policy
objective, yielding the \emph{DPO$^+$} variant reported in the headline results.

Critically, because the reweighting occurs at the data stage, every downstream
component — the RM, the DPO/PPO optimizer, and the preference-pair generator
in round $r{+}1$ — operates on an effectively unbiased signal from round one
onward.
This is the key structural difference from the baselines: rather than patching
the output of a biased RM, DPO$^+$ prevents the bias from entering the RM in
the first place.

\paragraph{Annealed correction.}
Across rounds, we anneal the correction strength following
\citet{kim2025cpm}:

\begin{equation}
    \alpha_r \;=\; \min\!\bigl(\alpha_0 + \Delta\alpha \cdot r,\; 1.0\bigr),
    \quad \alpha_0 = 0.4,\; \Delta\alpha = 0.12,
\label{eq:anneal}
\end{equation}

where $\alpha_r$ scales the IPW weights at round $r$.
The intuition is that early rounds have limited data and high variance in
propensity estimates; a smaller correction prevents overcorrection before the
propensity model has seen enough data to be reliable.
By round 5 the correction is near its maximum, matching the asymptotic regime
where the mechanism analysis in \S\ref{sec:theory} applies.
